% =========================================================================
% Preámbulo COMPARTIDO de los documentos independientes de riesgos
% (docs/thesis/plan-riesgos/ y docs/thesis/seguimiento-riesgos/).
%
% Estos documentos NO forman parte de main.tex (la memoria): son entregables
% sueltos en clase article, con el diseño heredado de los .tex de Gestión de
% Proyectos del autor, pero adaptados a la configuración global de la
% memoria (Carlito, márgenes 2,5 cm, parskip 6 pt, títulos de la guía).
% Un solo preámbulo para toda la familia: un ajuste de estilo aquí se
% propaga a todos los seguimientos futuros (uno por sprint hasta S37).
%
% RUTAS: los \input son relativos a la carpeta del documento, no a esta;
% compilar SIEMPRE desde la carpeta del .tex (plan-riesgos/ o
% seguimiento-riesgos/):
%   pdflatex --enable-installer -interaction=nonstopmode -file-line-error
%     -output-directory=build <documento>.tex        (2 pases; sin biber)
% =========================================================================

% --- Codificación (idéntico a packages.tex de la memoria) ---
\usepackage[utf8]{inputenc}
\usepackage[T1]{fontenc}

% --- Idioma: solo español (aquí no hay resumen trilingüe) ---
% es-noshorthands por coherencia con la memoria: acentos en UTF-8, sin
% caracteres activos (el conflicto con biblatex-apa no aplica aquí, pero
% mantener una única configuración de babel evita sorpresas al copiar texto
% entre la memoria y estos documentos).
\usepackage[spanish, es-tabla, es-noshorthands]{babel}

% --- Márgenes: 2,5 cm laterales e inferior, como la memoria ---
% El margen superior se reparte para alojar la cabecera enmarcada del
% diseño de referencia (2 cm de caja + separación) sin comerse el cuerpo:
% el texto arranca a 4 cm del borde y el ancho de texto sigue siendo 16 cm.
\usepackage[a4paper, left=2.5cm, right=2.5cm, bottom=2.5cm,
            top=1.2cm, includehead, headheight=2.2cm, headsep=0.6cm]{geometry}

% --- Tipografía: Carlito, la misma decisión que en style.tex ---
% (clon métrico libre de Calibri; pdflatex no puede usar la real).
% No se replica el arreglo de versalitas de style.tex: aquí no hay
% bibliografía APA y ningún texto usa \textsc.
\usepackage[sfdefault]{carlito}

% --- Espaciado: interlineado sencillo + 6 pt entre párrafos (guía) ---
\usepackage[skip=6pt]{parskip}

% --- Tablas ---
% longtable: el registro y el plan de respuesta superan una página.
% AVISO ESPERADO: "ignored error: Infinite glue shrinkage" en cada salto de
% página de una longtable (bug conocido e inofensivo, latex3/latex2e#1907;
% ver el comentario largo en packages.tex de la memoria).
\usepackage{array}
\usepackage{longtable}
\usepackage{tabularx}
% \RaggedRight (ragged2e) en vez de \raggedright en columnas muy estrechas:
% el \raggedright de LaTeX desactiva de facto la partición silábica y una
% palabra larga («Identificado») desborda una columna de 1,2 cm; la
% variante de ragged2e mantiene el guionado.
\usepackage{ragged2e}
% Las columnas X se componen como m{...}: el contenido de cada celda queda
% centrado verticalmente (aspecto de formulario del diseño de referencia).
\renewcommand{\tabularxcolumn}[1]{m{#1}}
% Tipos de columna del diseño de referencia:
%   L/C = ancho fijo alineado izquierda/centro; Y/Z = X alineada izq./centro.
% En columnas estrechas se prefiere Y (raggedright) a X justificada: la
% justificación en 3-4 cm abre ríos de blancos.
\newcolumntype{L}[1]{>{\raggedright\arraybackslash}m{#1}}
\newcolumntype{C}[1]{>{\centering\arraybackslash}m{#1}}
\newcolumntype{Y}{>{\raggedright\arraybackslash}X}
\newcolumntype{Z}{>{\centering\arraybackslash}X}
% Filas con aire en todas las tablas (el diseño de referencia lo fijaba
% tabla a tabla; aquí una sola vez).
\renewcommand{\arraystretch}{1.25}

% --- Utilidades ---
\usepackage{xcolor}    % \markfield (resaltado) y \todonote (macros.tex)
\usepackage{amssymb}   % \square, \blacksquare, \circlearrowright
\usepackage{xspace}    % lo necesita \bukit (macros.tex)
\usepackage{enumitem}
\setlist[itemize]{leftmargin=*, topsep=2pt, itemsep=2pt}
\setlist[enumerate]{leftmargin=*, topsep=2pt, itemsep=2pt}

% --- Títulos según la guía de estilo, como en style.tex de la memoria ---
% (sin \chapter: la clase article no lo tiene; el resto, idéntico).
\usepackage{titlesec}
\titleformat{\section}
  {\normalfont\bfseries\fontsize{12pt}{14.4pt}\selectfont\raggedright}
  {\thesection.}{0.5em}{\MakeUppercase}
\titleformat{\subsection}
  {\normalfont\bfseries\fontsize{12pt}{14.4pt}\selectfont}
  {\thesubsection.}{0.5em}{}
\titleformat{\subsubsection}
  {\normalfont\itshape\fontsize{12pt}{14.4pt}\selectfont}
  {\thesubsubsection.}{0.5em}{}

% --- Paginación dinámica (regla estricta del autor, 26-07-2026) ---
% PROHIBIDO el salto de página manual también en estos documentos: el
% espacio lo garantizan estos ganchos (o un \Needspace* local ante un bloque
% indivisible), nunca un \newpage. Los .tex de referencia sí usaban
% \newpage; es la única desviación deliberada respecto a ellos.
\usepackage{needspace}
\AddToHook{cmd/section/before}{\Needspace*{5\baselineskip}}
\AddToHook{cmd/subsection/before}{\Needspace*{5\baselineskip}}
\AddToHook{cmd/subsubsection/before}{\Needspace*{4\baselineskip}}
\widowpenalty=10000
\clubpenalty=10000
\displaywidowpenalty=10000
\raggedbottom
% Estos formularios componen mucho texto en columnas de 2-4 cm: se concede
% a TeX un margen de estirado de emergencia antes de desbordar la línea
% (evita Overfull \hbox menores sin tocar el contenido).
\setlength{\emergencystretch}{1em}

% --- Datos globales y macros de la memoria (única fuente: no duplicar) ---
% =========================================================================
% Variables globales del documento.
% Cambiar un dato AQUÍ lo actualiza en todo el documento (portada,
% metadatos del PDF...). No escribir nunca estos datos "a mano" en otros
% archivos: usar siempre estos comandos.
% =========================================================================

% --- Trabajo ---
\newcommand{\thesisTitle}{BUKIT: red social orientada a objetivos}% TODO: título definitivo
\newcommand{\thesisSubtitle}{Memoria del Trabajo de Fin de Grado}%
\newcommand{\thesisDegree}{Grado en Ingeniería Informática de Gestión y Sistemas de Información}%
\newcommand{\thesisCourse}{2026--2027}% curso académico
\newcommand{\thesisDate}{Bilbao, \today}% TODO: <Lugar, día, mes, año> de la entrega/defensa
\newcommand{\thesisYear}{2027}% año de elaboración y defensa

% --- Personas (la plantilla del centro usa "Apellido1 Apellido2, Nombre") ---
\newcommand{\thesisAuthor}{Liviu Deleanu}% nombre natural (para metadatos del PDF)
\newcommand{\thesisAuthorFormal}{Deleanu, Liviu}% TODO: confirmar formato con dos apellidos
\newcommand{\thesisDirector}{Villamañe Gironés, Mikel}%
\newcommand{\thesisCoDirector}{Apellido1 Apellido2, Nombre}% TODO: solo si hay codirector/a

% --- Universidad ---
\newcommand{\thesisUniversity}{Universidad del País Vasco / Euskal Herriko Unibertsitatea}%
\newcommand{\thesisFaculty}{Escuela de Ingeniería de Bilbao}%

% --- Palabras clave del resumen trilingüe ---
\newcommand{\thesisKeywordsEs}{red social, objetivos, monorepo, TypeScript, PostgreSQL}% TODO
\newcommand{\thesisKeywordsEu}{sare soziala, helburuak, monorepo, TypeScript, PostgreSQL}% TODO: revisar euskera
\newcommand{\thesisKeywordsEn}{social network, goals, monorepo, TypeScript, PostgreSQL}% TODO

% =========================================================================
% Comandos propios. Definir aquí cualquier cosa que se repita en el texto:
% si mañana cambia cómo se escribe, se cambia en un único sitio.
% =========================================================================

% Nombre del producto, siempre con el mismo formato.
\newcommand{\bukit}{\textsc{Bukit}\xspace}

% Términos en inglés, en cursiva de forma consistente: \eng{backend}
\newcommand{\eng}[1]{\textit{#1}}

% Notas de trabajo, visibles en rojo en el PDF. Buscar "todonote" y
% eliminarlas todas antes de entregar.
\newcommand{\todonote}[1]{\textcolor{red}{\textbf{[TODO: #1]}}}


% --- Cabecera y pie (diseño de referencia, sin logos) ---
% La caja de cabecera replica la del .tex de referencia (celda ancha de
% contexto + celda estrecha de identidad). La celda derecha llevaba el logo
% del centro: aquí va texto porque los logos UPV/EHU siguen pendientes de
% la fase de diseño (ver style.tex); cuando existan, sustituir por
% \includegraphics. 11 + 4,9 = 15,9 cm < 16 cm de texto: sin desbordes.
\usepackage{fancyhdr}
\newcommand{\riesgosHeader}{%
  \setlength{\tabcolsep}{0pt}%
  \renewcommand{\arraystretch}{1}%
  \begin{tabular}{|p{11cm}|p{4.9cm}|}
    \hline
    \parbox[c][2cm][c]{11cm}{\centering\bfseries
      {\fontsize{16}{19}\selectfont TRABAJO DE FIN DE GRADO\\[4pt]Curso \thesisCourse}%
    } &
    \parbox[c][2cm][c]{4.9cm}{\centering
      {\fontsize{14}{17}\selectfont\bfseries\bukit}\\[3pt]
      {\small\thesisFaculty}%
    }\\
    \hline
  \end{tabular}%
}
\pagestyle{fancy}
\fancyhf{}
\renewcommand{\headrulewidth}{0pt}
\fancyhead[C]{\riesgosHeader}
\fancyfoot[C]{\thepage}
% La primera página usa estilo "plain": se redefine para que también lleve
% la cabecera (mismo truco que el .tex de referencia).
\fancypagestyle{plain}{%
  \fancyhf{}%
  \renewcommand{\headrulewidth}{0pt}%
  \fancyhead[C]{\riesgosHeader}%
  \fancyfoot[C]{\thepage}%
}

% --- Símbolos y campos de formulario del diseño de referencia ---
\newcommand{\fieldline}[1]{\underline{\hspace{#1}}}
\newcommand{\bigline}{\noindent\rule{\textwidth}{0.4pt}}
\newcommand{\checkbox}{\(\square\)}
\newcommand{\checkedbox}{\(\blacksquare\)}
% Símbolo ⟳ = "si se rellena este apartado, nueva versión del plan".
% El original lo resolvía con un SVG externo (paquete svg + Inkscape);
% aquí basta el símbolo de amssymb: cero dependencias nuevas.
\newcommand{\updatesymbol}{\raisebox{-0.2ex}{\(\circlearrowright\)}}
% Resaltado amarillo del identificador de riesgo en los títulos.
\newcommand{\markfield}[1]{\colorbox{yellow!35}{\strut #1}}

% --- Enlaces y metadatos del PDF (al final, como exige hyperref) ---
% El pdftitle lo fija cada documento con \hypersetup tras cargar esto.
\usepackage[hidelinks]{hyperref}
\hypersetup{pdfauthor={\thesisAuthor}}
